% FILE: /home/jstvns/GPLC/MEMORIA/sections/futuro.tex

\newpage
\section{Limitaciones y Trabajo Futuro}\label{sec:limitaciones-trabajo-futuro}


Este capítulo tiene como objetivo presentar, de
manera crítica, las principales limitaciones
identificadas durante el desarrollo del intérprete
de GPLC y proponer líneas de trabajo futuro para
superar dichas limitaciones. En este trabajo, el
intérprete se concibió como un prototipo de
investigación: su foco principal es materializar
la semántica del lenguaje, permitir la ejecución
de programas y servir como plataforma para
experimentar con la gradualidad tanto en tipos
como en probabilidades. En consecuencia, existen
aspectos del sistema que se privilegiaron
(corrección semántica y fidelidad a la formulación
teórica) por sobre otros (optimización o escalabilidad a
programas de mayor tamaño).



En las siguientes subsecciones, se organizan las
limitaciones en categorías que reflejan el tipo de
dificultad observada: \textbf{i.} potenciales
imprecisiones semánticas o errores de
implementación, \textbf{ii.} evaluación de probabilidades
simbólicas en el output, \textbf{iii.} eficiencia de la resolución de
fórmulas, \textbf{iv.} rediseño de tipado para operaciones
binarias, \textbf{v.} ausencia de mecanismos de igualdad
para testing semántico, y \textbf{vi.} costos
computacionales asociados a \lstinline|Let| y
\lstinline|Dascr| en el AST destino.

\subsection{Imprecisiones Semánticas y/o Errores de Implementación}\label{subsec:lim-bugs}

Durante la implementación y la experimentación con
la ejecución de programas se identificaron
resultados no esperados. Estos casos pueden
deberse, por una parte, a errores de
implementación y, por otra, a errores en la
especificación original del lenguaje, los cuales
podrían manifestarse como comportamientos no
esperados en una implementación que sigue dicha
formulación de manera literal. Distinguir ambas
fuentes es esencial para definir líneas de
trabajo: en el primer caso corresponde corregir el
código y en el segundo, podría ser necesaria una
rectificación de la semántica
utilizada en la implementación.



Un ejemplo ilustrativo es el siguiente programa del código
\ref{lst:bug}, donde se observa que la distribución final contiene
valores que no coinciden con la intuición del usuario sobre la
ejecución. En el resultado, dos valores son compatibles con la
interpretación intuitiva: \lstinline|[Bool]true :: ?| y
\lstinline|[Real]13 :: Real|.  Sin embargo, aparecen además: \textbf{i.} un
error dinámico \lstinline|ERROR(?)|, que indica que una ascripción con
tipo desconocido no pudo validarse en tiempo de ejecución, y \textbf{ii.} un
valor \lstinline|[Real]13 :: ?| que indica que, en algún punto del
pipeline, el valor $13$ fue ascrito con tipo desconocido, a pesar
de que en el código \ref{lst:bug} no hay ningún elemento sintáctico
que haga la ascripción \lstinline|:: ?| necesaria o esperable.

\begin{gplcblock}[caption={Distribución con valores no esperados},label={lst:bug}]
? $ (true :: ?) : 13

>>>

{
  [Bool] true :: ? : SymProb19 ,
  ERROR(?): SymProb20 ,
  [Real] 13. :: ? : SymProb21 ,
  [Real] 13. :: Real : SymProb22
}
\end{gplcblock}

La evidencia hallada durante el desarrollo, indica
que este tipo de fenómeno se sitúa en la
evaluación de la ascripción de distribución en el
AST destino. Dicho nodo tiene la estructura
\lstinline|Dascr(e,m,nu)|, donde \lstinline|e| es
la evidencia, \lstinline|m| es la subexpresión, y
\lstinline|nu| es el tipo ascrito. Primero se
evalúa \verb|m| a una distribución de valores \verb|V|, y
luego se asigna a cada valor \verb|V| una evidencia y
un tipo de \verb|nu| para chequear consistencia en
tiempo de ejecución. En el proceso de asignar a cada valor de
\verb|V| una evidencia y un tipo de nu, hay asignaciones
que no se hacen correctamente, si los valores
\lstinline|[Real]13::?| y \lstinline|ERROR(?)| se van a considerar como no
esperados en V

\paragraph{Trabajo Futuro.}
En trabajo futuro, estos casos deben ser aislados y reproducidos de
manera sistemática, con el fin de establecer con claridad si el
comportamiento no esperado proviene de un error de implementación o de
un problema en la formulación teórica del lenguaje presentada en el
paper de referencia.  En particular, el caso del código~\ref{lst:bug}
requiere una revisión focalizada de la evaluación de \lstinline|Dascr|
en el AST destino, pues en dicha etapa se reconstruyen
correspondencias entre valores, tipos.
Una vez identificada la causa,
corresponde corregir el algoritmo (si es un bug) o 
rectificar la regla (si el fenómeno proviene de la
especificación). 

\subsection{Evaluación de Probabilidades Simbólicas}\label{subsec:lim-prob-simbolicas}

Actualmente, el intérprete retorna distribuciones simbólicas de
valores en las cuales las probabilidades están representadas por
variables simbólicas. Sin embargo, la fórmula asociada a la
distribución, que restringe los valores reales simultáneos que estas
variables pueden tomar, no se muestra en el output.  Si bien las
funcionalidades que realizan la impresión de fórmulas están
implementadas en el código, el volumen de información contenido en una
fórmula correspondiente a un programa no trivial, suele ser demasiado
grande para entregar una información interpretable por el usuario, y
por esta razón se omite en el output.  En consecuencia, el resultado
observado por el usuario corresponde esencialmente a listado de
valores asociados a un pesos variables, pero no a las restricciones
que hacen que el sistema sea soluble y semánticamente válido.

Esta decisión de diseño resulta útil en el contexto experimental del
prototipo, pues permite observar con claridad el los valores que
contiene el soporte de una distribución, y verificar que un programa
termina exitosamente. No obstante, incluso en este contexto, resulta
valioso contar con mecanismos que expresen el resultado de manera
numérica, instanciando las probabilidades simbólicas cuando el usuario
lo requiera.

 \paragraph{Trabajo Futuro.}
 Un trabajo futuro natural es mejorar la forma en que el intérprete
 comunica el resultado al usuario. En primer lugar, sería valioso
 permitir que la ejecución muestre la fórmula que restringe las
 probabilidades, de manera parcial y estratégica y en un formato que
 favorezca su legibilidad, pues la fórmula explica qué asignaciones
 reales son compatibles con la distribución retornada.  En segundo
 lugar, sería deseable habilitar un modo de instanciación de
 probabilidades simbólicas, ya sea extrayendo un modelo desde el
 solver cuando exista, o permitiendo que el usuario entregue valores
 para las probabilidades desconocidas \lstinline|?|.

\subsection{Resolución Incremental de Fórmulas}\label{subsec:lim-incremental}

En el intérprete, la satisfacibilidad de fórmulas se utiliza de forma
transversal: es un paso esencial para determinar consistencia,
precisión, el meet de tipos de distribución, y para validar
restricciones heredadas o acumuladas en tiempo de ejecución. En la
implementación actual, la función \lstinline|solve| que conecta el
intérprete con el solucionador Z3 se comporta como una \textit{caja
  negra} que recibe fórmulas y retorna únicamente un valor booleano
(true si existe solución, false si no). Este diseño favorece una
integración simple del solver dentro del prototipo, pero impide
reutilizar trabajo entre chequeos sucesivos: cada invocación construye
un contexto nuevo del solver, vuelve a traducir la fórmula completa a
expresiones nativas de Z3 y a resolverla, encareciendo el cómputo
cuando estas verificaciones se repiten frecuentemente durante la
evaluación, lo cual ocurre generalmente.

Esta limitación se vuelve especialmente relevante en operaciones entre
distribuciones (por ejemplo, durante un \lstinline|meet|), donde la
fórmula a resolver se construye por acumulación de restricciones pues
se incluyen las fórmulas de las distribuciones argumento y se agregan
nuevas variables y ecuaciones asociadas al coupling . En este
escenario, aunque las subfórmulas heredadas puedan asumirse satisfacibles
(un invariante garantizado por el flujo del intérprete), una solución para
una subfórmula no necesariamente se extiende a una solución del
sistema global, porque la operación introduce restricciones
adicionales.

\paragraph{Trabajo Futuro.}
Un trabajo futuro relevante es explorar un enfoque incremental que
permita reutilizar información entre chequeos sucesivos, reduciendo el
costo de resolver sistemas que se forman, en parte, por acumulación de
restricciones ya chequeadas. Sin embargo, esta mejora debe diseñarse
con cuidado, dado que, incluso si una subfórmula es soluble, su
solución no necesariamente es compatible con una solución del sistema
completo cuando se agregan nuevas restricciones. Por ello, una
dirección razonable es reutilizar contexto y restricciones de manera
segura, sin asumir extensibilidad automática de subfórmulas.

\subsection{Tipado de Operaciones Binarias}\label{subsec:lim-binops}

En la formulación implementada, las operaciones binarias aritméticas y
booleanas exigen que sus operandos sean consistentes con
\lstinline|Real| o \lstinline|Bool|, respectivamente. Una vez
satisfecha esta condición, el tipo de la operación se interpreta como
una distribución Dirac sobre el tipo correspondiente (intuitivamente,
una distribución cuyo soporte contiene un único tipo con probabilidad
1). Este diseño es coherente con la idea de que toda expresión denota
una distribución. Sin embargo, genera una tensión en el sistema:
existen construcciones en el lenguaje (como aplicación de abstracción,
ascripción simple, operaciones binarias, y condiciones de un
\lstinline|if|) que esperan un término de tipo simple (un valor), no
un término cuyo tipo inferido sea una distribución (aunque sea Dirac).


Como consecuencia, expresiones como \lstinline|(1+1)+7|,
\lstinline|(f (1 + 1))| o \lstinline|(1 + 1) :: Real| evidencian una
ambigüedad conceptual: la subexpresión \lstinline|(1 + 1)| puede verse
como un valor simple \lstinline|Real| o como una distribución
\lstinline|{Real:1}|. En el prototipo, esta tensión no se resolvió
mediante un rediseño del sistema de tipos, ya que el foco principal
del trabajo fue implementar la gradualidad en tipos de distribución y
en probabilidades desconocidas.

\paragraph{Trabajo Futuro.}
Esta limitación sugiere la necesidad de un rediseño que compatibilice
la interpretación de expresiones como distribuciones con los
contextos del lenguaje que esperan valores simples, como en los
ejemplos mencionados. Una línea de trabajo futuro consiste en
introducir un mecanismo explícito que permita tratar, por ejemplo,
\lstinline|{Real:1}| como \lstinline|Real| cuando el contexto lo
requiera, o bien adoptar un sistema de tipado más sensible al tipo
esperado por el contexto. En particular, uno de los desafíos está dado
por el hecho de que la elaboración transforma las operaciones binarias
en $\LetKW$ (el cual está típicamente asociado a la generación de
distribuciones), como se muestra a continuación:
\begin{equation*}
  v1 + v2  \mathrel{\leadsto} \LetTerm{x}{\eascr{\epsilon_1}{v_1}{\Real}}{\LetTerm{y}{\eascr{\epsilon_2}{v_2}{\Real}}{x+y}}
\end{equation*}
Sería importante, por lo tanto, distinguir
en el lenguaje destino entre los $\LetKW$ insertados
por la elaboración para chequear los valores
sumados, y aquellos que provienen de un $\LetKW$ en el
lenguaje fuente.


\subsection{Equivalencia Semántica y Testing}\label{subsec:lim-igualdad}

Una limitación importante para la validación sistemática es la
ausencia de un operador o mecanismo interno para comprobar igualdad
entre dos expresiones (o entre dos distribuciones resultantes). Sin
una noción de igualdad, resulta difícil automatizar tests que comparen
un resultado esperado con el que produce una expresión. Una
funcionalidad de este tipo tendría especial importancia para verificar
equivalencia semántica entre expresiones del lenguaje fuente,
comparando las distribuciones de valores que generan. Por ejemplo,
sería deseable poder contrastar directamente los programas:
\begin{itemize}
\item \lstinline|let x = (0.5 $ 1 : true) in 0.5 $ 1 : x| 
\item \lstinline|0.75 $ 1 : true|.
\end{itemize}


En el estado actual, lo más cercano a una comparación consiste en
ascribir manualmente un tipo de distribución esperado y verificar que
el programa termina sin error (código \ref{lst:pseudo-comparison}), lo
que es una condición necesaria pero no suficiente para equivalencia
semántica.

\begin{gplcblock}[caption={Verificación de tipo de distribución},label={lst:pseudo-comparison}]
(let x = (0.5 $ 1 : true) in  0.5 $ 1 : x ) :: {Real : 0.75, Bool : 0.25}

>>>

{
	 [Real] 1. :: Real : SymProb95 ,
	 [Real] 1. :: Real : SymProb96 ,
	 [Real] 1. :: Real : SymProb97 ,
	 [Bool] true :: Bool : SymProb98 
}  
\end{gplcblock}


\paragraph{Trabajo futuro.}
Una extensión importante del lenguaje sería
incorporar un mecanismo para comparar expresiones
del lenguaje fuente por equivalencia semántica, lo
que se traduce en comparar las distribuciones de
valores que dichas expresiones generan. Para ello
será necesario definir primero qué significa
igualdad entre valores (incluyendo el rol del tipo
ascrito) y decidir cómo tratar valores
funcionales, que se evalúan a clausuras. Con esa
base, podría definirse una noción de igualdad
entre distribuciones que considere tanto el
soporte como las restricciones sobre
probabilidades simbólicas. En presencia de
probabilidades desconocidas, también resulta natural estudiar
una igualdad parcial: distribuciones
que coinciden para algún conjunto no vacío de
instanciaciones de probabilidades simbólicas.

\subsection{Ineficiencia de \lstinline|Dascr| y \lstinline|Let|}\label{subsec:lim-eficiencia}

La evaluación de \lstinline|Let| y
\lstinline|Dascr| es un proceso computacionalmente
costoso.  Esto se debe a dos razones
estructurales. Primero, ambos constructos aparecen
con alta frecuencia en el AST destino: La
elaboración reestructura varios términos fuentes
en \verb|Let|, y inserta ascripciones de distribución.  Segundo, su evaluación ejecuta
operaciones entre distribuciones de tipo que, en
el peor caso, crecen combinatoriamente: si dos
distribuciones de tipo con tamaños $m$
y $n$, se combinan mediante una operación como
meet o reorder\footnote{explicar omisión}, la distribución resultante puede
alcanzar tamaño $m \times n$ antes de descartar
combinaciones de tipos simples inválidas.

La implementación de las operaciones entre
distribuciones de tipos se basa en la construcción
de un coupling simbólico (una matriz de variables
frescas, con dimensión $m\times n$) cuyas variables
se usan para construir y resolver el sistema de
restricciones que valida la operación (ver sección
\ref{par:symbolic}). Dicho sistema, en conjunción
con las fórmulas heredadas de las distribuciones
sobre las que se opera, pasa a ser la fórmula de la
distribución resultante.  Este enfoque es
directo y cuida la claridad y fidelidad a la
definición, pero es muy costoso en tiempo y
memoria cuando se repite de manera profunda y
recurrente en la recursión estructural de la
evaluación sobre el AST destino.

\paragraph{Trabajo futuro.}
Dado que \lstinline|Let| y \lstinline|Dascr| aparecen con alta
frecuencia en el AST destino y gatillan operaciones costosas entre
distribuciones, un trabajo futuro relevante es estudiar estrategias
para reducir el crecimiento del tamaño de distribuciones y fórmulas
durante la evaluación. En términos generales, esto apunta a evitar
combinaciones innecesarias en operaciones $m\times n$, y a simplificar
distribuciones y restricciones antes de delegar al solver. El objetivo
sería mejorar la escalabilidad del prototipo, permitiendo ejecutar
programas más grandes y observar comportamientos que hoy quedan fuera
por costo computacional.


\subsection{Resumen}\label{subsec:lim-resumen}

En síntesis, el intérprete desarrollado cumple su propósito como
plataforma experimental para ejecutar programas de GPLC y observar su
comportamiento con respecto a la gradualidad de probabilidades y
tipado. Sin embargo, el prototipo deja abiertas varias líneas de
mejora.  Por una parte, existen casos puntuales donde el
comportamiento relacionado con ascripciones requiere una
revisión cuidadosa para distinguir entre errores de implementación y
posibles imprecisiones en la formulación teórica original. Por otra,
la salida actual prioriza la visualización de valores y sus probabilidades simbólicas
, pero no incluye las restricciones (fórmula) que controlan
dichas probabilidades, lo que solo permite una interpretación limitada del
resultado cuando el usuario necesita probabilidades numéricas. Además,
el uso intensivo del solver hace que operaciones frecuentes
(especialmente, las asociadas a \lstinline|Let| y \lstinline|Dascr|
durante la evaluación) se vuelvan costosas y dificulten escalar a
programas más grandes. Finalmente, el diseño actual del tipado --donde
operaciones binarias producen por defecto distribuciones Dirac--
genera problemas en contextos que esperan valores simples, y la ausencia
de una noción de igualdad entre expresiones impide automatizar pruebas
de equivalencia semántica. El trabajo futuro debería apuntar a
cerrar estas brechas manteniendo la corrección semántica, pero
mejorando la informatividad del output, la reutilización de cómputo en
la resolución de fórmulas y la eficiencia general del intérprete.
